\documentclass{article}
\usepackage{graphicx}
\usepackage{url}

%% http://www.ctan.org/tex-archive/help/Catalogue/entries/igo.html
\usepackage{igo}
\gobansize{13}

\usepackage[empty,cm]{fullpage}
\begin{document}

%% The first page includes a 9x9 board.
\centerline{\Huge WPI Go Beginner's Night}
\centerline{\Huge 7pm Friday, February 27, 2009}
\centerline{\Huge @ Campus Center, 2nd floor}
\vfill

\begin{center}
\includegraphics{gb9.epsi}
\end{center}

\vfill

\newpage

\section*{Eye, eye, eye, eye, stayin' alive, stayin' alive...}

Go is often called the ``surrounding game'', and for good reason!  We
surround stones to capture them.  But another thing we
surround on the board is empty space.  Groups that surround a large
amount of empty space resist capture.  In fact, if we surround in a
particular way, our groups become altogether immune.
\begin{center}
\shortstack{
\cleargoban
\white{a3,b3,b2,b1,c3,d3,e3,e2,e1}
\showgoban
}
\hspace{.2in}
\shortstack{
\cleargoban
\white{a2,a3,b3,b1,c3,c2,c1}
\showgoban}
\hspace{.2in}
%
\shortstack{
\cleargoban
\white{c6,c5,c4,d6,d4,e5,e7,f5,f6,f7}
\showgoban
}
%
\hspace{.2in}
\shortstack{
\cleargoban
\white{a10,a11,b10,b12,c10,c11,c12,c13}
\showgoban
}
%% diagrams here
\end{center}
The distinguishing feature about these invincible groups is that
they've surrounded two distinct empty areas; we call these areas
\emph{eyes}.  Black can't capture because suicidal moves are
prohibited.  If a group can form two eyes, even despite the opponent's
best resistance, then that group is \emph{alive}.


\section*{Puzzles}
\begin{center}
\shortstack{
\cleargoban
\white{a3,b3,c3,d3,e3,e2,e1}
\black{a2,b2,c2,d2,d1}
\showgoban\\
Black's move.\\Where can Black form eyes?
}
\hspace{.3in}
\shortstack{
\cleargoban
\white{c12,c11,d11,e11,f11,g11,h11,j11,j12,j13}
\black{d12,e12,f12,f13,g12,h12,h13}
\showgoban\\
Black's move.\\Where can Black form eyes?
}
\hspace{.3in}
\shortstack{
\showgoban\\
White's move.\\Can White interfere with Black's eyes?\\
(\small{\it Your opponent's good move is your good move.})
}
\end{center}




\section*{Scoring}

Once the borders between Black and White are solid, how do we score a game?
\begin{enumerate}
\item Groups that can't avoid capture are given to the opponent.

\item Each player tallies the empty space they've completely surrounded
  by his or her groups.  This is \emph{territory}; we add territory
  and the number of captures together to get the total score.
\end{enumerate}
%
Here is an example of a game that's reached its end.
\begin{center}
\shortstack{
\cleargoban
\gobansize{9}
\black[1]{d7,f4,c4,f7,f6,g6,f5,g5,g7,f8}
\black[11]{e7,g8,e4,f3,e3,e2,d2,f2,c3,c9}
\black[21]{c8,e9,b9,d9,b7,e8,e1,f1,d1,d8}
\cleargobansymbols
\black[\igotriangle]{g7}
\showfullgoban\\
Both players agree that \blackstone[\igotriangle] is dead.}
\hspace{.3in}
\shortstack{
\clear{g7}
\showfullgoban\\
\blackstone[\igotriangle] is given to white.}
\hspace{.3in}
\shortstack{
\gobansymbol{a1,a2,a3,a4,a5,a6,a7,a8,a9,b1,b2,b3,b4,b5,b6,b8,c1,c2,c5,c6,c7,d3,d4,d5,d6,e5,e6}{b}
\gobansymbol{f9,g1,g2,g3,g4,g7,g9,h1,h2,h3,h4,h5,h6,h7,h8,h9,j1,j2,j3,j4,j5,j6,j7,j8,j9}{w}
\showfullgoban\\
Black and White count territories.}
\end{center}
Black has 27 territories, and his score is 27.  White has 25
territories, and since White received 1 captured stone, her total
score is 26.

\section*{For more information...}
\begin{itemize}
\item \emph{WPI Go} (\url{http://go.hashcollision.org/}) is our homepage.

\item \emph{Sensei's Library} (\url{http://senseis.xmp.net/}), a wiki
  devoted to the game.

\item \emph{The Interactive Way to Go}
  (\url{http://playgo.to/interactive/}), which gives a more leisurely
  introduction.

\item \emph{Go Problems}
 (\url{http://www.goproblems.com}) is a database of puzzles, many
  which involve the life and death of groups.

\end{itemize}
\end{document}
